%Formattazione e codifica caratteri
\usepackage[english]{babel}
\usepackage[utf8]{inputenc}
\usepackage{microtype}
\usepackage{sansmathfonts}
% Tavola dei contenuti
\usepackage{tocloft}
\renewcommand{\cftdot}{}
\setlength{\cftbeforechapskip}{5pt}
% Figure
\usepackage{graphicx}
\usepackage{caption}
\usepackage{subfig}
\usepackage{svg}
\usepackage{float}
% Matematica
\usepackage{amsmath,amsfonts,amssymb,amsthm}
\usepackage{color}
% Tabelle
\renewcommand*\familydefault{\sfdefault}
\usepackage{booktabs}
%Listati
\usepackage{listings}
\lstloadlanguages{MATLAB}
% Testatine
%\author{A. Siniscalchi}
\usepackage{fancyhdr}
\usepackage{epigraph}
\usepackage{titlesec, blindtext}
\usepackage[dvipsnames]{xcolor}
\usepackage{tikz-network}
\usepackage{pgfplots}
 \pgfplotsset{compat=newest}
  \usetikzlibrary{plotmarks}
  \usetikzlibrary{arrows.meta}
  \usepgfplotslibrary{patchplots}
  \usepackage{grffile}

\usepackage{tikz}
\usepackage{tikz-cd}
\usepackage{hyperref}
\pgfplotsset{plot coordinates/math parser=false}
\pagestyle{fancy}
\setlength{\headheight}{15pt}
\fancyhf{}
\fancyhead[LE,RO]{}
\fancyhead[RE,LO]{\leftmark}
\fancyfoot[CE,CO]{}
\fancyfoot[LE,RO]{\thepage}

\renewcommand{\headrulewidth}{1.5pt}
% Titoli dei capitoli che "rubano" meno spazio

\definecolor{gray75}{gray}{0.75}
\newcommand{\hsp}{\hspace{15pt}}

\SetVertexStyle[MinSize=1\DefaultUnit,FillColor=white]


% Per importare le figure in MATLAB
  \newlength\figureheight
 \newlength\figurewidth
%Per le osservazioni
\theoremstyle{plain}
\newtheorem{thm}{Theorem}[section]
\newtheorem{lem}[thm]{Lemma} 
\newtheorem{obs}[thm]{Observation} 
\newtheorem{prop}[thm]{Proposition}
\newtheorem{cor}[thm]{Corollary} 
\newtheorem{rem}[thm]{Remark}
\newtheorem{ex}[thm]{Example} 
\newtheorem{cla}[thm]{Claim} 
\newtheorem{as}[thm]{Assertion}
\newtheorem{note}[thm]{Note}
\newtheorem{defn}[thm]{Definition}

%COMANDI PERSONALIZZATI 
\newcommand{\Q}{\mathbb{Q}}
\newcommand{\R}{\mathbb{R}}
\newcommand{\B}{\mathbb{B}}
\newcommand{\C}{\mathbb{C}}
\newcommand{\Z}{\mathbb{Z}}
\newcommand{\D}{\mathbb{D}}
\newcommand{\s}{\mathbb{S}}
\newcommand{\torus}{T_{2, 2g+1}}
\newcommand{\tonde}[1]{\left( #1 \right)}
\newcommand{\ses}{\Leftrightarrow} 
\newcommand{\abs}[1]{\left\vert#1\right\vert}
\newcommand{\angol}[1]{\langle #1 \rangle}
\newcommand{\spa}{\mathbin{\textcolor{white}{-}}} 
\newcommand{\q}[1]{``#1''}
\newcommand{\inter}{\nu(K_m)^{\mathrm{o}}}
\newcommand{\interd}{\nu(D')^{\mathrm{o}}}
\newcommand{\estd}{\mathbb{B}^4 \setminus \nu(D)^{\mathrm{o}}}
\newcommand{\estdi}{\mathbb{B}^4 \setminus \nu(D')^{\mathrm{o}}}
\newcommand\restr[2]{{% we make the whole thing an ordinary symbol
  \left.\kern-\nulldelimiterspace % automatically resize the bar with \right
  #1 % the function
  \littletaller % pretend it's a little taller at normal size
  \right|_{#2} % this is the delimiter
  }}
\newcommand{\littletaller}{\mathchoice{\vphantom{\big|}}{}{}{}}

% per mettere le figurelle
%\begin{figure}[h]
    %\centering
    %\includegraphics[scale=0.25]{figures/prova.jpeg}
    %\caption{Caption}
    %\label{fig:evviva ho messo una figura}
%\end{figure}
%faccio riferimento \ref{fig:evviva ho messo una figura}